% Generated draft: requires mathematical and source review.
\documentclass[11pt]{article}
\usepackage[margin=25mm]{geometry}
\usepackage{fontspec}
\setmainfont{Latin Modern Roman}
\usepackage{amsmath,amssymb,mathtools,unicode-math}
\setmathfont{Latin Modern Math}
\usepackage{enumitem,needspace,xurl,hyperref}
\setlength{\parindent}{0pt}
\setlength{\parskip}{5pt}
\setlength{\emergencystretch}{3em}
\newcommand{\N}{\mathbb{N}}
\newcommand{\Z}{\mathbb{Z}}
\newcommand{\Q}{\mathbb{Q}}
\newcommand{\R}{\mathbb{R}}
\newcommand{\C}{\mathbb{C}}
\newcommand{\F}{\mathbb{F}}
\newcommand{\A}{\mathbb{A}}
\newcommand{\PP}{\mathbb{P}}
\newcommand{\E}{\mathbb{E}}
\newcommand{\eps}{\varepsilon}
\newcommand{\dd}{\,\mathrm{d}}
\begin{document}
\section*{UnsolvedMath ID: 800689}
\textbf{Source problem code:} catalogue entry
\subsection{Definitions and mathematical statement}
Let \(\mathbb{P}^{k}=\mathbb{P}^{k}(\C)\) denote complex projective space. For a holomorphic foliation \(\mathscr{F}\), a regular leaf is a connected maximal complex integral manifold of its tangent distribution on its regular locus. A compact curve in a leaf is understood here to mean a compact complex analytic subset of that leaf of complex dimension one.

The construction problem is to exhibit a non-elementary holomorphic foliation \(\mathscr{F}\) on some \(\mathbb{P}^{k}\), with leaf dimension \(d>1\), such that
\[
\forall L\in\operatorname{Leaves}(\mathscr{F}),\qquad
\nexists\, C\subseteq L\text{ with }C\text{ a compact complex analytic curve}.
\]
The ambient dimension \(k\) and leaf dimension \(d\) are not specified in the supplied statement; necessarily \(2\leq d\leq k\). The source does not explicitly require examples for every admissible dimension.

The qualification ``non-elementary'' is essential but is not defined in the supplied record. It is retained as a source qualification rather than replaced by an invented condition. The supplied record also does not specify permissible singularities, how singular leaves are treated, or whether the foliation must have positive codimension. Thus the display records the requested avoidance property, but the full class of admissible foliations requires clarification from the source. The interpretation of ``compact curve'' above also requires confirmation.
\subsection{Short English statement}
Construct a non-elementary holomorphic foliation on a complex projective space whose leaves have complex dimension greater than one and none of whose leaves contains a compact complex curve. The supplied record leaves ``non-elementary,'' the dimension range, and the singularity conventions unspecified.
\subsection{Sources}
\begin{itemize}
\item Supplied local UnsolvedMath registry record: permanent ID \(800689\), exact source problem code \texttt{catalogue entry}, title ``Holomorphic Foliation Theory --- Problem 16,'' in the AMR Open Problem Lists collection.
\item Dinh and Sibony, \emph{Some Open Problems on Holomorphic Foliation Theory} (2018), Problem 16. Supplied source URL: \url{https://arxiv.org/abs/1807.05640}. The registry attributes the problem to Dinh and Sibony, with proposed year 2018.
\item The supplied provenance records access on 2026-10-16 and extraction by \texttt{arxiv\_latex}. These are supplied provenance claims; the source has not been independently consulted for this draft.
\item The registry labels the problem ``open.'' This is a supplied status label, not an independently verified assessment of the literature.
\end{itemize}
\end{document}
